\chapter{同余子群：怎样从有限余数回到精确数值}\label{ch:congruence}
\section{先看一个可以算到最后的输送}
给定两个本原整数列
\[v=\binom23,qquad w=\binom78.
\]
它们模五相同，但实际相差很大。矩阵
\[\gamma=\begin{pmatrix}-4&5\\-5&6\end{pmatrix}
\]
满足 $\det\gamma=1$、$\gamma\equiv I\pmod5$，并且 $\gamma v=w$。

这正是我们想要的工具：不改变指定余数，却能把实际向量精确送到另一个。定义
\[\Gamma(q,R)=\{M\in\SL_2(R):M\equiv I\pmod{qR}\}.
\]
$\SL_2(R)$ 是行列式为一的二阶矩阵群；“主同余子群”就是其中模 $q$ 等于单位矩阵的那些。这个名称虽然长，条件只是四个矩阵项的同余。

若 $R=\Z[1/D]$ 且 $\gcd(q,D)=1$，环里的分数怎样取模？把 $a/D^k$ 送到 $a(D^k)^{-1}\pmod q$。所以 $R/qR$ 与普通的 $\Z/q\Z$ 是同一个有限环。若 $q$ 含已经在 $R$ 中可逆的因子，先把它们去掉即可。例如在 $\Z[1/2]$ 中，模六的条件等价于模三，因为二已经是单位。

\section{为什么必须先取得这个群的真实实现}
模 $q$ 相同的两列有无穷多个。如果它们都能由 $\Gamma(q,R)$ 互相送到，而这个群的每个矩阵又都有真实平均词，那么无限问题就只剩有限问题：先送对余数，剩下由群精确完成。

不过，光在计算机里看到模 $q$ 的矩阵很多，并不能推出这个群已经可实现。模运算会把一大批不同有理矩阵压到一起，丢掉的部分正是我们需要控制的。这是两个外部定理真正进入证明的地方。

“有限指数”表示一个子群把整个群分成有限多个陪集，可以先把它理解成只剩有限种余类。Morris 的定理从足够多的初等矩阵得出有限指数；Serre 的定理再说，对我们的环，这样的子群一定包含某个主同余子群。
\InsertProof{thm:arithmetic-external}
第一项使用 Morris\cite[Theorem 6.1(2)]{morris} 的二维无限单位情形。第二项使用 Serre\cite{serre} 的有理 $S$ 整数结论，准确陈述见\cite[pp.308--309]{prasad}。这里 $S$ 含实处和 $D$ 的素因子处，$D>1$ 保证至少有一个有限处。若 $D=1$，变成 $\SL_2(\Z)$，第二项一般就不成立。

这两项是本讲引用的深结果。后面从“不知道多大的模数”改进到明确的平方模数，则是我们自己能逐步算出的消元。

\section{为什么最后出现的是平方模数}
我们已经有上、下两族步长在 $hR$ 内的剪切。外部定理先保证一个未知 $\Gamma(q_0,R)$ 已包含其中。接下来想证明：任何 $g\equiv I\pmod{h^2}$ 也属于已生成的群。

取更深的模数 $q=h^2q_0$。在有限环 $R/qR$ 里，先给第一行加第二行的一个合适倍数，使左上角变成单位。然后像高斯消元一样消掉非对角项。剩下的对角项也能写成四个剪切，因为它与一的差被 $h^2$ 整除。

消元只保证与 $g$ 模 $q$ 相同。但剩余误差属于一开始已经取得的 $\Gamma(q_0,R)$，所以能够精确补掉。这就是“先有一个真群，再用有限运算改善它”的顺序。
\InsertProof{lem:deep}

\section{如果起初只有上剪切怎么办}
上剪切始终固定第一坐标轴。如果一个可逆返回 $J$ 把这条轴送到另一条直线，那么 $JU(t)J^{-1}$ 就固定那条新的直线。把两条直线当成坐标轴，就看到了上、下两种方向。

只是这个换基矩阵未必有整数逆，所以要把分母损失写进同余模数。下面的 $g^3$ 就是为这件事付出的明确代价。它不是多出三个物理位置，也不是三次平均的次数。
\InsertProof{lem:transverse}

\section{相同余数的本原列，为什么真的能精确互达}
在 $R$ 中，“列 $(a,b)$ 本原”表示可以找到 $u,v\in R$ 使 $ua+vb=1$。在整数环这就是 $\gcd(a,b)=1$。它的线性代数含义是：这列能够补成一个行列式为一的二阶矩阵。

给起点 $v$ 补一列 $x$，给目标 $w$ 补一列 $y$。若能让两个补成的矩阵也模 $q$ 相同，那么
\[(w,y)(v,x)^{-1}
\]
便既把 $v$ 送到 $w$，又模 $q$ 等于 $I$。下面证明的全部工作，就是把第二列稍作修改，以满足这两个要求。
\InsertProof{08_congruence-unit4}
\InsertProof{lem:fiber}
这是整套方法里最值得停下来理解的一步：有限同余计算负责把你送进正确的“房间”，这个引理负责在房间内送到那个\emph{确切的点}。最终是等号，不是取极限。

\section{取得主同余子群以后，逆更容易造了}
在有限矩阵群里，某矩阵的正整数次幂一定回到 $I$。如果 $Z^m\equiv I\pmod q$，那么真实 $Z^m$ 就落进已能实现的 $\Gamma(q,R)$，其逆也已可用。因此
\[Z^{-1}=Z^{m-1}(Z^m)^{-1}.
\]
这把“有限商中的周期”变成了一条真正的逆词。

行列式不是一时，要额外处理整数内容。Smith 标准形是通过整数可逆行列变换把矩阵化成对角形的定理；这里用它设计一个两边夹乘，使结果能除掉行列式而仍是整数矩阵。再由 CRT 让这个夹乘同时模 $q$ 正确。
\InsertProof{lem:inverse}
第一次读到这儿，可先记住它的用途和条件：\textbf{已经取得真实主同余子群后，行列式与控制模数互素的整数返回，可以获得正向射影逆。}具体夹乘等式留在上面的完整证明中，第二遍逐项相乘核对即可。
